% s-032-p032-m002; original PDF 32, 33
\Aphorism{s-032-p032-m002}{21}
Alle Dinge haben begrenzte Lebensdauer, sei es active oder passive; Wahrheiten, Werthschätzungen, Ideale können ebensowenig in infinitum gelebt werden als ein Organismus in infinitum leben kann. Dieser erstaunliche Satz scheint so neu zu sein, so alt er klingt: nichts steht für das Empfinden früherer Culturen fester als die Ewigkeit des Menschen und seiner Lebensformen: seiner Seele, seiner Sittlichkeit, seiner Welt und seines Gottes. Zu dieser unbedingten Gefühlsvoraussetzung stimmt aufs beste die leidenschaftliche Unersättlichkeit, mit der die ältere Menschheit an ihren Werthgebilden und Idealen hängt und sie bis aufs Blut aussaugt: man hält sie eben für unerschöpflich und ahnt nicht, dass, was gegenwärtig davon weggelebt wird, ausgegebenes Kapital und für spätere Generationen verloren ist. Dem christlichen Gott merkt man es heute an, dass sich zwei Jahrtausende im Guten und Schlimmen an ihm genährt haben: solche Zeiträume mögen wohl noch kräftigere Ackerkrume erschöpfen. Man wirft uns Modernen Unbescheidenheit vor: eine Religion, die so viele gesättigt, werde wohl auch für uns noch gut genug sein. Aber eben weil wir Frucht- und Fleischesser und keine Koprophagen sind, verzichten wir auf den Abfall bereits verzehrter Mahlzeiten; eben weil Millionen in diesem Becken gebadet haben, steigen wir nicht mit hinein. Und wie mit Religion, steht es fast mit allen Lebensformen, in die wir heute hineingeboren werden; sie alle sind, nicht etwa durch die fortschreitende Wissenschaft, sondern einfach durch die fortschreitende Zeit für uns entwerthet; das, was an ihnen Leben war, haben frühere Zeiten verschwenderisch genug aufgezehrt und abgenagt. Es hilft nichts, wir müssen uns schon unsere eigene Nahrung suchen und unseren eigenen Acker bestellen, eigene Wahrheiten wissen, eigene Götter glauben, eigene Leidenschaften grossziehen; es steht uns gar nicht mehr frei, in epigonenhafter Unfruchtbarkeit vom ererbten Gute zu zehren. Wir sind arm, arm wie die Kirchenmäuse, unbegütert, ungespeist, ungetränkt: wovon leben wir eigentlich? Sind wir selbst nicht unser eigenes Futter? Und wer will es uns wehren, wenn wir aus der Noth eine Tugend machen, wenn wir uns selber gut schmecken, uns selber \Sperr{lieben}?
\par

